% --- Volume III: Analysis / Linear Algebra ---
% Section: Vector Space Preliminaries

\begin{remark*}[Toolkit --- Vector Space Preliminaries]
\textbf{Concept family:} The concrete foundations for $\mathbf{F}^n$ before the
abstract vector space axioms are introduced.

\medskip
\textbf{Atomic items in this section:}
\begin{itemize}
  \item Definition: List
  \item Definition: $\mathbf{F}^n$
  \item Definition: Addition in $\mathbf{F}^n$
  \item Theorem: Addition Is Commutative in $\mathbf{F}^n$
  \item Definition: Zero Vector in $\mathbf{F}^n$
  \item Definition: Additive Inverse in $\mathbf{F}^n$
  \item Definition: Scalar Multiplication in $\mathbf{F}^n$
\end{itemize}

\medskip
\textbf{Key predicate:}
\[
  \operatorname{ListEquality}((x_1,\ldots,x_n),(y_1,\ldots,y_m))
  \;\colonequiv\;
  n = m \;\wedge\; \forall k \le n\;(x_k = y_k).
\]
\end{remark*}

% ---------------------------------------------------------

\begin{definitionbox}{Definition (List)}
\begin{definition}[List]
\label{def:list}
Suppose $n$ is a non-negative integer. A \textbf{list} of length $n$ is an ordered collection of $n$ elements.
\end{definition}
\end{definitionbox}

\begin{remark*}[Standard quantified statement]
\[
\operatorname{List}(L) \iff \exists n \in \mathbb{N} \cup \{0\} : \bigl( \operatorname{Length}(L) = n \land L = (x_1, x_2, \dots, x_n) \bigr)
\]
\end{remark*}

\begin{remark*}[Definition predicate reading]
\[
\operatorname{List}(L) \iff \exists n \in \mathbb{N}_0, \exists (x_1, \dots, x_n) : L = (x_1, \dots, x_n)
\]
\end{remark*}

\begin{remark*}[Interpretation]
A list is the primary mechanism for discussing ordered finite sets of vectors or scalars. Unlike a set, the structural integrity of a list depends on both the identity and the position of its members.
\begin{itemize}
    \item \textbf{Order Sensitivity}: The lists $(a, b)$ and $(b, a)$ are distinct unless $a=b$.
    \item \textbf{Finiteness}: By Axler's convention, a list must have a finite number of elements. Infinite ordered collections are categorized as sequences.
    \item \textbf{Empty List}: A list of length $n=0$ is called the empty list, denoted by $()$.
\end{itemize}
\end{remark*}

\begin{dependencies}
\begin{itemize}
  \item \hyperref[def:function]{Function}
  \item \hyperref[def:whole-numbers]{Whole Numbers}
\end{itemize}
\end{dependencies}
% ---------------------------------------------------------

\begin{definitionbox}{Definition ($\mathbf{F}^n$)}
\begin{definition}[$\mathbf{F}^n$]
\label{def:fn}
Let $\mathbf{F}$ denote $\mathbb{R}$ or $\mathbb{C}$, and let $n$ be a positive integer.
Then $\mathbf{F}^n$ is defined to be the set of all lists of length $n$ of elements of $\mathbf{F}$:
\[
\mathbf{F}^n
=
\{(x_1,\dots,x_n) : x_1,\dots,x_n \in \mathbf{F}\}.
\]
The entries in a list in $\mathbf{F}^n$ are called the coordinates of the list.
\end{definition}
\end{definitionbox}

\begin{remark*}[Standard quantified statement]
\[
\begin{aligned}
&x \in \mathbf{F}^n \\
&\quad\Longleftrightarrow\quad
\exists x_1,\dots,x_n \in \mathbf{F}
\text{ such that }
x = (x_1,\dots,x_n).
\end{aligned}
\]
\end{remark*}

\begin{remark*}[Definition predicate reading]
\[
\operatorname{Fn}(\mathbf{F},n)
\;\colonequiv\;
\mathbf{F}^n = \{(x_1,\ldots,x_n) : x_k \in \mathbf{F}\}.
\]
\end{remark*}

\begin{remark*}[Interpretation]
An element of $\mathbf{F}^n$ is an ordered list of $n$ scalars from $\mathbf{F}$.
The order matters, and the list is completely specified by its coordinates.
\end{remark*}

\begin{dependencies}
\begin{itemize} 
  \item \hyperref[def:list]{Definition: List}
\end{itemize}
\end{dependencies}

% ---------------------------------------------------------

\begin{definitionbox}{Definition (Addition in $\mathbf{F}^n$)}
\begin{definition}[Addition in $\mathbf{F}^n$]
\label{def:addition-in-fn}
Let
\[
x=(x_1,\dots,x_n), \qquad y=(y_1,\dots,y_n)
\]
be vectors in $\mathbf{F}^n$.
The sum of $x$ and $y$ is defined by
\[
x+y
=
(x_1+y_1,\dots,x_n+y_n).
\]
\end{definition}
\end{definitionbox}

\begin{remark*}[Standard quantified statement]
\[
\begin{aligned}
&\forall x,y \in \mathbf{F}^n, \\
&\quad
x=(x_1,\dots,x_n)\land y=(y_1,\dots,y_n)
\;\Rightarrow\;
x+y=(x_1+y_1,\dots,x_n+y_n).
\end{aligned}
\]
\end{remark*}

\begin{remark*}[Interpretation]
Addition in $\mathbf{F}^n$ is defined coordinatewise: add the first coordinates,
then the second coordinates, and so on through the $n$th coordinates.
\end{remark*}

\begin{dependencies}
\begin{itemize}
  \item \hyperref[def:fn]{Definition: $\mathbf{F}^n$}
\end{itemize}
\end{dependencies}

% ---------------------------------------------------------

\begin{theorembox}{Theorem (Addition Is Commutative in $\mathbf{F}^n$)}
\begin{theorem}[Addition Is Commutative in $\mathbf{F}^n$]
\label{thm:addition-commutative-in-fn}
For all $x,y \in \mathbf{F}^n$,
\[
x+y=y+x.
\]
\hyperref[prf:addition-commutative-in-fn]{Go to proof.}
\end{theorem}
\end{theorembox}

\begin{remark*}[Standard quantified statement]
\[
\forall x,y \in \mathbf{F}^n,\; x+y=y+x.
\]
\end{remark*}

\begin{remark*}[Interpretation]
Because addition in $\mathbf{F}^n$ is coordinatewise and scalar addition in
$\mathbf{F}$ is commutative, vector addition in $\mathbf{F}^n$ is commutative.
\end{remark*}

\begin{dependencies}
\begin{itemize}
  \item \hyperref[def:addition-in-fn]{Definition: Addition in $\mathbf{F}^n$}
\end{itemize}
\end{dependencies}

% ---------------------------------------------------------

\begin{definitionbox}{Definition (Zero Vector in $\mathbf{F}^n$)}
\begin{definition}[Zero Vector in $\mathbf{F}^n$]
\label{def:zero-vector-in-fn}
The symbol $0$ denotes the list of length $n$ all of whose coordinates are $0$:
\[
0=(0,\dots,0)\in \mathbf{F}^n.
\]
\end{definition}
\end{definitionbox}

\begin{remark*}[Standard quantified statement]
\[
0=(0,\dots,0)\in \mathbf{F}^n.
\]
\end{remark*}

\begin{remark*}[Interpretation]
The zero vector is the vector whose every coordinate is the scalar $0$ in
$\mathbf{F}$. It is the additive identity for $\mathbf{F}^n$.
\end{remark*}

\begin{dependencies}
\begin{itemize}
  \item \hyperref[def:fn]{Definition: $\mathbf{F}^n$}
\end{itemize}
\end{dependencies}

% ---------------------------------------------------------

\begin{definitionbox}{Definition (Additive Inverse in $\mathbf{F}^n$)}
\begin{definition}[Additive Inverse in $\mathbf{F}^n$]
\label{def:additive-inverse-in-fn}
Let
\[
x=(x_1,\dots,x_n)\in \mathbf{F}^n.
\]
The additive inverse of $x$ is defined by
\[
-x=(-x_1,\dots,-x_n).
\]
\end{definition}
\end{definitionbox}

\begin{remark*}[Standard quantified statement]
\[
\begin{aligned}
&\forall x \in \mathbf{F}^n, \\
&\quad
x=(x_1,\dots,x_n)
\;\Rightarrow\;
-x=(-x_1,\dots,-x_n).
\end{aligned}
\]
\end{remark*}

\begin{remark*}[Interpretation]
The additive inverse in $\mathbf{F}^n$ is formed coordinatewise by taking
the additive inverse of each coordinate in $\mathbf{F}$.
\end{remark*}

\begin{dependencies}
\begin{itemize}
  \item \hyperref[def:fn]{Definition: $\mathbf{F}^n$}
  \item \hyperref[def:zero-vector-in-fn]{Definition: Zero Vector in $\mathbf{F}^n$}
\end{itemize}
\end{dependencies}

% ---------------------------------------------------------

\begin{definitionbox}{Definition (Scalar Multiplication in $\mathbf{F}^n$)}
\begin{definition}[Scalar Multiplication in $\mathbf{F}^n$]
\label{def:scalar-multiplication-in-fn}
Let
\[
\lambda \in \mathbf{F}
\quad\text{and}\quad
x=(x_1,\dots,x_n)\in \mathbf{F}^n.
\]
The product of $\lambda$ and $x$ is defined by
\[
\lambda x
=
(\lambda x_1,\dots,\lambda x_n).
\]
\end{definition}
\end{definitionbox}

\begin{remark*}[Standard quantified statement]
\[
\begin{aligned}
&\forall \lambda \in \mathbf{F}\;\forall x \in \mathbf{F}^n, \\
&\quad
x=(x_1,\dots,x_n)
\;\Rightarrow\;
\lambda x=(\lambda x_1,\dots,\lambda x_n).
\end{aligned}
\]
\end{remark*}

\begin{remark*}[Interpretation]
Scalar multiplication in $\mathbf{F}^n$ is defined coordinatewise: multiply
each coordinate of the vector by the scalar \(\lambda\).
\end{remark*}

\begin{dependencies}
\begin{itemize}
  \item \hyperref[def:fn]{Definition: $\mathbf{F}^n$}
\end{itemize}
\end{dependencies}
